\clearpage
\section{Extended Data}
\renewcommand{\figurename}{Extended Data Fig.}
\setcounter{figure}{0}

\noindent Nine Extended Data figures are planned (the comparison with the source recipes moved to Supplementary Table~4 when that column was cut back to the runs already finished). Three of them need no images and are drawn below from the same data table as the main figures; the others are listed with their content and data source.

\begin{center}\small\sffamily
\begin{tabular}{@{}l l l l@{}}
\toprule
\textbf{Extended Data Fig.} & \textbf{Content} & \textbf{Source} & \textbf{Status} \\
\midrule
1 & Example images and labels of each source; annotation style & prepared datasets; label audit & needs images \\
2 & Field-of-view masks and canvas scale per source & FOV reports & needs images \\
3 & Leakage audit: near-duplicates between sources and splits & hash and feature audit & audit not run \\
4 & Sensitivity to the metric implementation (thresholds, FOV, overlap) & evaluation switches & with Fig.~\ref{fig:literature}d,e \\
5 & Per-image distributions and threshold sensitivity & per-image histograms & after E2 \\
6 & (withdrawn; source recipes are Supplementary Table~4) & E2b & -- \\
7 & Learning-rate sweep (DDR) & E2c & drawn \\
8 & Loss with the architecture fixed (U-Net) & E9 & drawn \\
9 & DiaRetDB1 official protocol; ultra-widefield at two scales & E3a & after E3a \\
10 & Computational cost & E10 & drawn \\
\bottomrule
\end{tabular}
\end{center}

\setcounter{figure}{6}
\begin{figure}[h!]
\centering
\includegraphics[width=\textwidth]{\figfile{ed07_lr_sweep}}
\caption{\textbf{Learning-rate sweep.} Validation \mAUPR{} on DDR of one general (HRNet) and one fundus-specific (M2MRF) network with the learning rate of the recipe multiplied by 1/3, 1 and 3. Kendall $\tau$ compares the order of the models at the recipe value with their order when each uses its best of the three values.}
\end{figure}

\setcounter{figure}{7}
\begin{figure}[h!]
\centering
\includegraphics[width=\textwidth]{\figfile{ed08_loss}}
\caption{\textbf{Loss with the architecture fixed.} \textbf{a}, \mAUPR{} of U-Net on the DDR test set with cross-entropy, batch Dice, cross-entropy with batch Dice (the main-table recipe) and cross-entropy with a per-image Dice term (the registered recipe, which collapses on rare classes); seed 0. \textbf{b}, Per-class change in AUPR against cross-entropy.}
\end{figure}

\setcounter{figure}{9}
\begin{figure}[h!]
\centering
\includegraphics[width=\textwidth]{\figfile{ed10_cost}}
\caption{\textbf{Computational cost.} \textbf{a}, \mAUPR{} on the IDRiD test set against the training time of one seed on IDRiD (40\,GB A100, one task per card); marker area, number of parameters. \textbf{b}, Parameters, trainable parameters under the unified protocol, floating-point operations for one 1,440-pixel canvas, peak training memory and inference time per image.}
\end{figure}

\clearpage
\section{Supplementary information}
\noindent\textbf{Supplementary tables.}
\begin{enumerate}[leftmargin=*,label=\textbf{\arabic*}.,itemsep=1pt]
\item Datasets: source, camera, field of view, resolution, annotation procedure, splits, licence; DiaRetDB1 positive-image counts against the original paper.
\item Models: all \rn{n.models} registered variants with tier, code, weights, reproducibility level, tracks and environment.
\item Reported numbers extracted from the literature (one row per paper, baseline, dataset, class and metric), with split, resolution, metric implementation and provenance.
\item Recipe and deviations of every run; source recipes (B1) against the unified protocol for the runs that completed.
\item Per-class results of every model on every test set: AUPR with bootstrap interval, 11-threshold AUPR, Dice and intersection over union at both thresholds, lesion-level precision, recall and F1 by lesion size.
\item Split lists.
\item Computational cost.
\item Refined IDRiD: the crop geometry recovered for each of the 27 test images, the label differences per class, and every IDRiD-trained model's per-class AUPR against the refined and the official labels through the same geometry.
\end{enumerate}
\noindent\textbf{Supplementary notes.} One note per model (source, reproducibility level, environment, changes, differences from the paper); the budget probe with every version of the rule; protocol details.
